\chapter{Aftrekken: eerste stappen}\label{ch:g1:subtraction}

Zeven ballonnen, twee vliegen weg: hoeveel blijven er over? Wegnemen is
\emph{aftrekken}, met het teken $-$. En aftrekken heeft nog een tweede
gezicht: het antwoordt ook op ``hoeveel meer?''.

\section{Wegnemen}

\begin{definition}[Aftrekken]\label{def:g1:subtraction:def}
\emph{Aftrekken} is wegnemen en tellen wat overblijft. We schrijven
\[
7 - 2 = 5
\]
en lezen: ``zeven min twee is vijf''.
\end{definition}

\begin{omfigure}
\begin{tikzpicture}[scale=0.6]
  \foreach \i in {0,...,4} \fill[omDef] (\i*1.05,0) circle (0.3);
  \foreach \i in {5,6} {
    \fill[omDef!30] (\i*1.05,0) circle (0.3);
    \draw[omThm, thick] (\i*1.05-0.32,-0.32) -- (\i*1.05+0.32,0.32);
    \draw[omThm, thick] (\i*1.05-0.32,0.32) -- (\i*1.05+0.32,-0.32);
  }
  \node[right, font=\small] at (7.6,0)
    {$7 - 2 = 5$: streep er twee door, vijf blijven over};
\end{tikzpicture}

{\small Een aftrekking in plaatjes: doorstrepen is wegnemen.}
\end{omfigure}

\begin{example}[Terugtellen]\label{ex:g1:subtraction:countback}
Voor een kleine aftrekking tel je terug op de lijn: $11 - 3$: zeg
``elf'' en doe drie stappen terug --- ``tien, negen, acht''. Dus
$11 - 3 = 8$.
\end{example}

\section{Hoeveel meer?}

\begin{example}[Het verschil]\label{ex:g1:subtraction:difference}
``Nina heeft $9$ knikkers en Tom heeft er $6$. Hoeveel heeft Nina er
meer?'' Er wordt niets weggenomen, en toch is het een aftrekking: maak
paren van de knikkers en tel de knikkers van Nina die zonder partner
blijven:
\[
9 - 6 = 3 .
\]
Nina heeft er $3$ meer. De uitkomst van een aftrekking heet het
\emph{verschil}\index{verschil}.
\end{example}

\begin{method}[Doortellen]\label{met:g1:subtraction:countup}
Om een \omterm{ex:g1:subtraction:difference}{verschil} te vinden, is \emph{door}tellen vaak het makkelijkst:
voor $12 - 9$ begin je bij $9$ en klim je naar $12$: ``tien, elf,
twaalf'' --- drie stappen. Dus $12 - 9 = 3$. (Klim bij grotere gaten
via $10$: $14 - 8$: van $8$ naar $10$ is $2$, van $10$ naar $14$ is
$4$: \omterm{ex:g1:subtraction:difference}{verschil} $6$.)
\end{method}

\begin{omfigure}
\begin{tikzpicture}[scale=0.6]
  \draw[-Stealth] (5.6,0) -- (15.4,0);
  \foreach \i in {6,...,15} \draw (\i,-0.1) -- (\i,0.1);
  \foreach \i in {8,10,14}
    \node[below, font=\footnotesize] at (\i,-0.12) {\i};
  \draw[-Stealth, omDef, thick] (8,0.3) .. controls (9,0.9) ..
    (10,0.3);
  \node[omDef, font=\small] at (9,1.0) {$+2$};
  \draw[-Stealth, omThm, thick] (10,0.3) .. controls (12,1.2) ..
    (14,0.3);
  \node[omThm, font=\small] at (12,1.25) {$+4$};
\end{tikzpicture}

{\small $14 - 8$ met doortellen: van $8$ naar $14$ zijn het
$2 + 4 = 6$ stappen.}
\end{omfigure}

\section{Optellen en aftrekken zijn tweelingen}

\begin{example}[Rekenfamilies]\label{ex:g1:subtraction:family}
De getallen $3$, $4$ en $7$ vormen een klein gezin:
\[
3 + 4 = 7, \qquad
4 + 3 = 7, \qquad
7 - 4 = 3, \qquad
7 - 3 = 4 .
\]
Wie één \omterm{def:g1:addition:def}{som} van het gezin kent, krijgt de drie andere er gratis bij.
Aftrekken maakt optellen ongedaan: na $+4$ brengt $-4$ je terug.
\end{example}

\begin{example}[Controleren]\label{ex:g1:subtraction:check}
Heb je $13 - 5 = 8$ uitgerekend? Controleer met de tweelingoptelling:
$8 + 5 = 13$. \checkmark\ Een aftrekking is precies goed als het
terugtellen klopt.
\end{example}

\begin{example}[Welke bewerking?]\label{ex:g1:subtraction:choose}
``$6$ vogels op de draad, $4$ vliegen weg'': wegnemen, $6 - 4 = 2$.
``$6$ vogels, en er komen $4$ bij'': samenvoegen, $6 + 4 = 10$.
``$6$ spreeuwen en $4$ roodborstjes: hoeveel spreeuwen zijn er meer?'':
\omterm{ex:g1:subtraction:difference}{verschil}, $6 - 4 = 2$. Teken het verhaaltje als je twijfelt!
\end{example}

\section{Oefeningen}

\begin{exercise}[$\star$]\label{exo:g1:subtraction:1}
Teken, streep door en reken uit: $6 - 2$; \quad $8 - 5$; \quad
$10 - 4$.
\end{exercise}

\begin{exercise}[$\star$]\label{exo:g1:subtraction:2}
Reken uit door terug te tellen: $9 - 3$; \quad $12 - 2$; \quad
$15 - 4$.
\end{exercise}

\begin{exercise}[$\star$]\label{exo:g1:subtraction:3}
Zoek het \omterm{ex:g1:subtraction:difference}{verschil} door door te tellen: $11 - 8$; \quad $13 - 9$; \quad
$16 - 12$.
\end{exercise}

\begin{exercise}[$\star$]\label{exo:g1:subtraction:4}
Schrijf de vier sommen van het gezin van $2$, $6$ en $8$.
\end{exercise}

\begin{exercise}[$\star$]\label{exo:g1:subtraction:5}
Reken uit en controleer met de tweelingoptelling: $14 - 6$; \quad
$17 - 9$.
\end{exercise}

\begin{exercise}[$\star$]\label{exo:g1:subtraction:6}
Vul de gaten in:
\[
10 - \;?\; = 7, \qquad
\;?\; - 5 = 5, \qquad
12 - \;?\; = 12 .
\]
\end{exercise}

\begin{exercise}[$\star$]\label{exo:g1:subtraction:7}
``Er liggen $15$ kersen in de schaal. De kinderen eten er $8$ op.
Hoeveel blijven er over?'' Schrijf de aftrekking en de antwoordzin.
\end{exercise}

\begin{exercise}[$\star$]\label{exo:g1:subtraction:8}
Kies bij elk verhaaltje $+$ of $-$ en reken uit: ``$9$ eenden op de
vijver, $5$ zwemmen weg''; ``$9$ eenden, er komen $5$ gansjes bij'';
``Ali is $9$, zijn zusje is $5$: hoeveel jaar ouder is Ali?''.
\end{exercise}

\begin{exercise}[$\star\star$]\label{exo:g1:subtraction:9}
Moedereend had $12$ eieren. Een paar zijn uitgekomen; nu blijven er
$12 - ?$ eieren en ze telt $7$ niet-uitgekomen eieren. Hoeveel eieren
zijn uitgekomen? (Welke rekenfamilie is dit?)
\end{exercise}
